% !PRINTED 645
\weberpage{645}
6. \gesperrt{If two integral functionals are associated with a third, then they are also associated with one another.}

7. \gesperrt{The norm $N(\alpha)$ is divisible by $\alpha$.}

This follows from equation (6), \S\,178: $\varphi(\alpha)=0$, whose last coefficient is $\overline{A}_n=\pm N(\alpha)$. For according to this:
\[
\pm N(\alpha) = \alpha(\alpha^{n-1}+\overline{A}_1\alpha^{n-2}+\cdots+\overline{A}_{n-1}).
\]

8. \gesperrt{An integral functional whose absolute norm is a constant is a unit, and conversely the absolute norm of a unit $\varepsilon$ is a constant different from zero.}

For if $\varepsilon$ is a unit, then
\[
N(\varepsilon) \text{ and } N\left(\frac{1}{\varepsilon}\right) = \frac{1}{N(\varepsilon)}
\]
are integral functionals in $\overline{Z}$; consequently $N(\varepsilon)$ is a unit in $\overline{Z}$. Conversely $N(\varepsilon)$ is divisible by $\varepsilon$; hence, if $N(\varepsilon)$ is a unit in $\overline{Z}$, $\varepsilon$ is a divisor of 1, that is, a unit.

9. \gesperrt{There are integral rational functions of $z$, for example the absolute norm of $\alpha$, that are divisible by $\alpha$. If $a$ is an integral rational function of $z$ divisible by $\alpha$ and of the lowest possible degree, then every other integral rational function of $z$ divisible by $\alpha$ is divisible by $a$.}

\begin{center}
\S\,181. \gesperrt{Greatest common divisor.}
\end{center}

If $\alpha, \beta, \ldots$ are integral functionals in $\overline{\Omega}$ and $x, y, \ldots$ are variables that do not occur in $\alpha, \beta, \ldots$, then
\weq{1}{\delta = \alpha x + \beta y + \cdots}
is, by \S\,180, 2., likewise an integral functional, and indeed $\delta$ is \gesperrt{divisible by every common divisor of $\alpha, \beta, \ldots$} If $x_0, y_0, \ldots$ are arbitrary integral functions or functionals in $\overline{Z}$, then, as is now to be proved,
\weq{2}{\delta_0 = \alpha x_0 + \beta y_0 + \cdots}
is divisible by $\delta$. For if we denote the absolute norm of $\delta$ by $D$, then
\weq{3}{N(\delta) = D E(x, y, \ldots) = D E,}
and $E(x, y, \ldots)$ is a unit in $\overline{Z}$, but at the same time an integral homogeneous function of the variables $x, y, \ldots$
